% Source: https://artofproblemsolving.com/community/c3h1956427_the_biggest_number_smaller_than_10
% Stable identity: aops-topic-1956427-q01
Suppose $x<10$ were the largest real number less than $10$. Then $y=\dfrac{x+10}{2}$ satisfies $x<y<10$ (it is the midpoint of $x$ and $10$), so $y$ is a real number less than $10$ and larger than $x$, a contradiction. Hence no largest real number below $10$ exists. (Numbers such as $9.9$, $9.99$, $9.999,\dots$ get as close to $10$ as we like.)

Among integers the situation is different: the integers less than $10$ are $\dots,8,9$, and the largest is $\boxed{9}$.
